% RQ4 Abstract: Treasury Resilience

DAO treasuries hold billions of dollars in assets, yet treasury management practices remain ad hoc. Market volatility, proposal spending, and revenue fluctuations create fiscal risks that threaten organizational sustainability.

We investigate treasury resilience through multi-agent simulation, examining how management parameters affect fiscal stability. Through \experimentruns{} simulation runs, we analyze three policy dimensions: (1) buffer fractions that reserve portions of treasury for emergencies, (2) emergency top-up mechanisms that restore depleted buffers, and (3) spending limits that cap proposal-driven outflows.

Our results characterize the tradeoff between resilience and capital efficiency. Large buffers reduce volatility but idle capital that could fund operations. Aggressive spending limits protect reserves but may block valuable proposals. Emergency top-ups provide backstop protection but require inflow capacity.

We identify optimal parameter ranges for different risk tolerances and provide practical guidelines for treasury policy design. Our findings inform DAOs seeking to balance operational flexibility with long-term fiscal sustainability.
